% BHU PYQ -- Commerce: Business Entrepreneurship (2025-26 NEP)
\documentclass[11pt,a4paper]{article}
\usepackage[a4paper,top=2.5cm,bottom=2.5cm,left=2.8cm,right=2.8cm,headheight=14pt]{geometry}
\usepackage{amsmath,amssymb}
\usepackage{fontspec}
\setmainfont{Kohinoor Devanagari}
\usepackage{microtype,fancyhdr,enumitem,hyperref,lastpage,parskip}

\hypersetup{colorlinks=true,urlcolor=black,linkcolor=black,pdftitle={COMSE11: Business Entrepreneurship -- BHU NEP 2025-26}}

\pagestyle{fancy}\fancyhf{}
\renewcommand{\headrulewidth}{0.4pt}\renewcommand{\footrulewidth}{0.4pt}
\fancyhead[L]{\small \textit{Banaras Hindu University}}
\fancyhead[R]{\small \textit{COMSE11: Business Entrepreneurship (2025-26)}}
\fancyfoot[C]{\small Page \thepage\ of \pageref{LastPage}}

\newlist{parts}{enumerate}{2}
\setlist[parts,1]{label=(\alph*),leftmargin=*,itemsep=4pt,topsep=3pt}
\setlist[parts,2]{label=(\roman*),leftmargin=*,itemsep=2pt,topsep=2pt}

\newcommand{\pts}[1]{\hfill \textit{[#1]}}

\begin{document}

\begin{center}
  {\large\bfseries BANARAS HINDU UNIVERSITY}\\[2pt]
  {\normalsize B. Com. (Hons.) Semester I Examination, 2025-26 (NEP)}\\[6pt]
  \rule{0.8\linewidth}{0.5pt}\\[6pt]
  {\large\bfseries COMMERCE}\\[4pt]
  {\large\bfseries Paper Code: COMSE11 : Business Entrepreneurship}\\[4pt]
  {\normalsize Time: 2:30 Hours \hfill Full Marks: 70}\\[2pt]
  {\small REF NO. CECONF/2025-26/PS/26}
\end{center}
\rule{\linewidth}{0.4pt}
\smallskip

\noindent \textit{Instructions: Answer \textbf{all} questions. The questions are of equal value (14 marks each). Write your Roll No.\ at the top immediately on receipt of this question paper.}

\smallskip
\noindent \textit{सभी प्रश्नों के उत्तर दीजिए। सभी प्रश्नों के अंक समान हैं।}

\bigskip

\begin{enumerate}[label=\textbf{\arabic*.},leftmargin=*,itemsep=14pt]

  \item Explain the concept of entrepreneurship and also bring out the various types of entrepreneurs.\pts{14}
  
  उद्यमिता की अवधारणा की व्याख्या कीजिए और इसके विभिन्न प्रकारों को भी स्पष्ट कीजिए।

  \begin{center}\textbf{OR / अथवा}\end{center}

  Describe the importance of entrepreneurship in a capitalist economy.\pts{14}
  
  पूंजीवादी अर्थव्यवस्था में उद्यमिता के महत्व का वर्णन कीजिए।

  \item Explain the concept of change keeping in mind the profitable survival of a business.\pts{14}
  
  व्यवसाय के लाभदायक अस्तित्व को ध्यान में रखते हुए परिवर्तन की अवधारणा की व्याख्या कीजिए।

  \begin{center}\textbf{OR / अथवा}\end{center}

  Write a critical note on social responsibility of entrepreneurs.\pts{14}
  
  उद्यमियों के सामाजिक दायित्व पर एक आलोचनात्मक टिप्पणी लिखिए।

  \item How is a new business idea identified and developed?\pts{14}
  
  एक नया व्यावसायिक विचार किस प्रकार पहचाना और विकसित किया जाता है?

  \begin{center}\textbf{OR / अथवा}\end{center}

  Write a detailed note on entrepreneurial decision process.\pts{14}
  
  उद्यमीय निर्णय प्रक्रिया पर एक विस्तृत टिप्पणी लिखिए।

  \item Write a detailed note on financial planning.\pts{14}
  
  वित्तीय नियोजन पर एक विस्तृत टिप्पणी लिखिए।

  \begin{center}\textbf{OR / अथवा}\end{center}

  What is business growth? How is growth possibility identified?\pts{14}
  
  व्यावसायिक वृद्धि क्या है? वृद्धि की संभावना कैसे पहचानी जाती है?

  \item Answer the following questions in not more than 75 words each:\pts{4 $\times$ 3.5 = 14}
  
  निम्नलिखित प्रश्नों के उत्तर दीजिए, प्रत्येक उत्तर अधिकतम 75 शब्दों में:
  \begin{parts}
    \item List the various functions of entrepreneurs.
    
    उद्यमियों के विभिन्न कार्यों की सूची बनाइए।

    \item Do you think that entrepreneurs are essentially born entrepreneurs? Explain.
    
    क्या आप सोचते हैं कि उद्यमी अनिवार्यतः जन्मजात उद्यमी होते हैं? समझाइए।

    \item Explain in brief any one technique of business idea generation.
    
    व्यावसायिक विचार उत्पत्ति की किसी एक तकनीक को संक्षेप में समझाइए।

    \item Write about the sources of finance.
    
    वित्त के स्रोतों के बारे में लिखिए।
  \end{parts}

\end{enumerate}

\end{document}
